\documentclass[11pt,a4paper]{article}
\usepackage[margin=25mm]{geometry}
\usepackage[T1]{fontenc}
\usepackage{lmodern,microtype,amsmath,amssymb,amsthm,booktabs,array,tabularx}
\usepackage[hidelinks]{hyperref}
\usepackage{url}
\newtheorem{theorem}{Theorem}
\newtheorem{lemma}[theorem]{Lemma}
\newtheorem{proposition}[theorem]{Proposition}
\newtheorem{observation}[theorem]{Observation}
\newcommand{\Rp}{\mathbb R_{>0}}
\newcommand{\cN}{\mathcal N}
\newcommand{\dd}{\operatorname{diag}}
\newcolumntype{Y}{>{\raggedright\arraybackslash}X}
\setlength{\emergencystretch}{1.5em}
\title{A certified three-equilibrium bound for 375 fixed\protect\\
quadratic zero-one reaction networks}
\author{Tatsuhiko Takimoto}
\date{24 October 2026 \quad Draft v1}
\begin{document}
\maketitle
\begin{abstract}
Consider the fixed list of 375 three-species, six-reaction networks
published by Tang and Zhang. For each network in this list and every
strictly positive choice of rate constants, we give a computer-assisted
proof that there are at most three distinct positive equilibria,
including degenerate equilibria. Network-specific elimination
certificates yield 272 cubic and 103 quartic residual polynomials.
The certificates exclude vanishing elimination denominators at positive
equilibria and parameter specializations that would make the residual
identically zero;
the quartics have opposite signs at their leading and constant
coefficients. An independently implemented resultant calculation
and a separate Sylvester-determinant check agree for all 375 networks.
Within this list, absence of inflow together with acyclicity of the
unimolecular production graph selects precisely the two known
exceptions to bistability; this is a post-selected finite observation.
We also treat one seven-reaction extension, proving at most one
exponentially stable positive equilibrium for all positive rates and
certifying a rational example with three positive equilibria.
The original enumeration and the stability classification of the
other 373 networks are outside our independent certification.
\end{abstract}

\section{The finite question and its scope}
Tang and Zhang~\cite{TZ} identify 375 quadratic zero-one networks
admitting three positive equilibria, of which 373 admit bistability.
The classification and the identification of the two exceptions,
33449 and 239089, are theirs. Their Theorem~2 states the upper bound
three; Section~5.1 gives an elimination argument for an example, while
the Discussion asks for a theoretical explanation of the bound.
We address the finite list they supplied, with the rate constants
left arbitrary and positive. We are grateful for their classification
and data, which make this question accessible.

The contribution is a checked algebraic reduction for every member
of that list, covering special positive parameter values as well as
generic ones. We do not give a new enumeration of networks or a new
general elimination method. Univariate reduction and the relation
between the derivative of a reduced equation and the Jacobian
determinant belong to established approaches, including
Torres--Feliu~\cite{TF} and Banaji--Feliu~\cite{BF}.
The latter's bounds for specified numbers of distinct reactants do
not immediately give three for the present three-species,
six-distinct-reactant inputs. Our assertion is a finite collection
of all-parameter statements, rather than a test of 375 rate vectors.

Let $Y,N\in\mathbb Z^{3\times6}$ be the reactant and stoichiometric
matrices. Column $j$ describes the reaction from $Y_{\cdot j}$ to
$Y_{\cdot j}+N_{\cdot j}$ with rate $\kappa_j>0$. The mass-action
equations are
\begin{equation}\label{eq:mass}
 \dot x_i=F_i(x;\kappa)=\sum_{j=1}^6 N_{ij}\kappa_j
                  \prod_{\ell=1}^3 x_\ell^{Y_{\ell j}}.
\end{equation}
Both complexes have entries zero or one, and every reactant has
molecularity at most two. Here ``quadratic'' means at most two,
not that every reaction is bimolecular. All 375 stoichiometric
matrices $N$ have rank three. A positive equilibrium has $x\in\Rp^3$ and $F(x;\kappa)=0$;
it is nondegenerate when $\det J_F\ne0$, and exponentially stable
when all eigenvalues of $J_F$ have negative real part.

Write $\cN$ for the 375 records in the fixed upstream file identified
in Appendix~\ref{app:input}. Network numbers below are the upstream
\texttt{candidate\_index}, not row numbers or matrix ranks.

\begin{theorem}\label{thm:bound}
For every network $G\in\cN$ and every $\kappa\in\Rp^6$, the
mass-action system~\eqref{eq:mass} has at most three distinct
positive equilibria. No nondegeneracy assumption is required.
\end{theorem}

Sections~\ref{sec:cert} and \ref{sec:independent} explain the two
computational routes supporting this theorem. The structural
observation and the seven-reaction proposition are separate results;
neither is needed for the upper bound.

\section{A univariate certificate for each network}\label{sec:cert}
Fix a network and retain one concentration $t$, writing the other
two as $u,v$. The equations are affine in each of $u,v$ separately,
with possible terms $1,u,v,uv$. A recorded invertible integer row
transformation $G=TF$ cancels $uv$ in two equations. When $uv$ is
absent already, two suitable equations are selected directly.
The certificates give
\begin{equation}\label{eq:cramer}
 A(t)\binom{u}{v}=b(t),\qquad
 D(t)=\det A(t),\qquad
 D(t)u=U(t),\quad D(t)v=V(t).
\end{equation}
All entries are polynomials in $t$ and the rate constants. The
record includes $T$, its nonzero determinant, the coordinate order,
$D,U,V$, and the third transformed equation $H$.

\subsection{Excluding division by zero on the positive branch}
In 251 networks, $D$ has nonzero coefficients of one strict sign.
In the remaining networks this certificate holds for $U$ (100
networks) or $V$ (24 networks). A nonzero polynomial whose nonzero
coefficients have one sign cannot vanish when all its variables
are positive.

If a solution of~\eqref{eq:cramer} had $D=0$, the adjugate
identities would give $U=V=0$. Thus a sign certificate for either
numerator also excludes $D=0$ at every positive equilibrium.
This step covers special positive rate choices; it is not merely
division in a field of rational functions. Every positive equilibrium
therefore satisfies
\[
 u=U/D,\qquad v=V/D,
\]
so its retained coordinate $t$ determines it uniquely.

Substitution in $H$ gives the polynomial
\begin{equation}\label{eq:residual}
 P(t)=D(t)^2 H\bigl(U(t)/D(t),V(t)/D(t),t\bigr)
      =c L(t)Q(t),\qquad c\in\mathbb Q\setminus\{0\}.
\end{equation}
Every removed factor in $L$ either has coefficients of one strict
sign or divides $D$. Such factors are nonzero at a positive
equilibrium. The records contain 706 strict-sign factor entries
and 21 divisor-of-$D$ entries, with multiplicities recorded separately.
No parameter factor that might vanish at positive rates is discarded
without a nonvanishing argument.

Consequently, positive equilibria inject into the positive roots of
$Q$. Conversely, a root of $Q$ with
\begin{equation}\label{eq:admissible}
 t>0,\qquad D(t)\ne0,\qquad U(t)/D(t)>0,\quad V(t)/D(t)>0
\end{equation}
gives a positive equilibrium: the first two transformed equations
hold, $P=0$ gives $H=0$, and $T$ is invertible. Roots failing these
conditions are not counted as positive equilibria.

\subsection{Root bounds that survive specialization}
The finite checks give the following two cases.

\begin{table}[ht]
\centering\small
\begin{tabularx}{\linewidth}{@{}rYY@{}}
\toprule
Networks & Residual certificate & Consequence at every positive rate vector\\
\midrule
272 & Degree at most three; at least one coefficient has a strict
sign & $Q$ is nonzero and has at most three positive roots.\\
103 & Degree four; leading coefficient strictly negative and
constant coefficient strictly positive & At least one negative
real root, hence at most three positive roots.\\
\bottomrule
\end{tabularx}
\caption{Root bounds from the Cramer certificates. Roots are bounded
even with multiplicity; the theorem counts distinct equilibria.}
\label{tab:degrees}
\end{table}

For a quartic, $Q(0)>0$ and $Q(t)\to-\infty$ as $t\to-\infty$,
so the negative root follows by continuity. For a cubic certificate,
a leading coefficient can vanish at particular positive rates,
but that only lowers the degree. A recorded strict-sign coefficient
prevents the entire polynomial from becoming zero. Thus the
argument includes degree drops and degenerate equilibria.

\begin{proof}[Proof of Theorem~\ref{thm:bound}]
For each of the 375 inputs, exact polynomial checks verify the
transformation and Cramer identities, the denominator certificate,
the factorization~\eqref{eq:residual}, the nonvanishing of removed
factors at positive equilibria, and the appropriate row of
Table~\ref{tab:degrees}. Equations~\eqref{eq:cramer}--\eqref{eq:residual}
give an injection from positive equilibria to positive roots of
a nonzero polynomial with at most three such roots. This proves
the assertion for every positive rate vector of each input network.
\end{proof}

The finite calculations here are symbolic. A SymPy construction
produces the certificates, and a separately written sparse-polynomial
checker using rational arithmetic verifies them without importing
SymPy. These two programs were written within one construction
effort. The next section describes the independent verification.

\section{Independent elimination and certificate checks}\label{sec:independent}
An independent implementation starts again from the fixed $Y,N$
matrices. It uses sequential elimination and resultants rather
than the two-variable Cramer construction. For an ordered set of
coordinates $u,v,w$, it writes a first equation as
\[
 F_u=-A(v,w)u+B(v,w),\qquad A(v,w)>0.
\]
It substitutes $u=B/A$ into the remaining equations and clears
the positive denominator, obtaining $H(v,w),K(v,w)$.
For the selected equation
\[
 H(v,w)=a(w)v^2+b(w)v+c(w),
\]
the coefficient certificates give $a\leq0$, $c\geq0$, with at
least one inequality strict throughout the positive domain.
If both are strict, the two roots have opposite signs. If $c=0$
there is at most one positive root of the remaining linear factor;
if $a=0$ the nonzero constant term prevents an identically zero
linear equation. Hence, for fixed $w>0$, there is at most one
positive $v$, and then at most one $u$.

Every positive equilibrium supplies a root of
$R(w)=\operatorname{Res}_v(H,K)$. Only factors certified nonzero
on the positive domain are removed. Each retained factor has a coefficient whose nonzero monomial
coefficients all share one strict sign. This coefficient cannot
vanish at positive rates, so the surviving polynomial in $w$ is
not identically zero at any positive rate specialization. The
records store these tests as \texttt{coefficient\_signs} within
\texttt{root\_factors}; Appendix~\ref{app:input} locates them. Descartes' rule is
applied allowing every possible sign, including zero, for each
coefficient not assigned a strict sign. The maximum possible
number of sign changes bounds all positive roots, including
multiple ones. Extra roots introduced by a resultant can only
weaken this upper bound; the proved uniqueness of the positive
fiber is what permits a bound on equilibria.

This route covers 372 inputs directly. Three require an invertible
integer change of equations before the same elimination steps:
\begin{center}\small
\begin{tabular}{@{}rlll@{}}
\toprule
Network & First equation & Eliminate; retain & Positive pivot $A$\\
\midrule
48328 & $F_1+F_3$ & $x_1,x_2$; $x_3$ & $2\kappa_5x_3$\\
74814 & $F_1+F_3$ & $x_1,x_2$; $x_3$ & $\kappa_5x_3$\\
100108 & $F_1+F_2$ & $x_1,x_3$; $x_2$ & $\kappa_6x_2$\\
\bottomrule
\end{tabular}
\end{center}
Each transformation has determinant one. The result is again
272 cubic and 103 quartic residuals, all with positive-root bound
three. The retained coordinate differs from the Cramer route in
79 inputs, while the degree and conclusion agree in every input.

A second checker in this independent verification reconstructs
the substituted equations directly from the original matrices.
It expands each Sylvester determinant by permutation enumeration,
with sizes at most $4\times4$, and verifies the factorization and
coefficient-sign bounds. It calls neither the first checker's
resultant/factorization functions nor its helper functions.
The two independent-verification checkers still share SymPy's
polynomial arithmetic; this is not a comparison of two unrelated
computer-algebra systems.

\section{A structural observation about the known exceptions}
Call a reaction an \emph{inflow} when its reactant complex is zero.
Define the unimolecular production graph on the three species:
place $i\to j$, for $i\ne j$, when a reaction whose reactant is
exactly $X_i$ produces $X_j$. This graph uses only unimolecular
reactions, not the signs of the full Jacobian.

\begin{observation}\label{obs:struct}
Within $\cN$, the conjunction of no inflow and an acyclic
unimolecular production graph selects exactly networks 33449
and 239089.
\end{observation}

The direct census, independently checked by two cycle-detection
procedures, is
\begin{center}
\begin{tabular}{@{}lrrr@{}}
\toprule
 & Acyclic & Cyclic & Total\\
\midrule
Inflow present & 255 & 112 & 367\\
No inflow & 2 & 6 & 8\\
\bottomrule
\end{tabular}
\end{center}
The condition was selected after a twelve-network pilot and
inspection of the eight no-inflow networks. Its match within
this list is therefore not an independent predictive test.
Observation~\ref{obs:struct} is not a theorem that this graph
condition determines stability for arbitrary reaction networks.
In particular, the graph omits bimolecular reactions.

For the two selected networks, a separate certificate collection
\cite{Adata} gives complete positive-branch reductions: each has
at most three distinct positive equilibria and at most one
exponentially stable positive equilibrium at arbitrary positive
rates. If three distinct positive equilibria exist, exactly one
is exponentially stable. We recall these results briefly to
identify the known exceptions. The collection is prepared for a
separate data-and-tools release; its final DOI is not yet available.

The determinant criterion of Tang--Zhang, Lemma~3~\cite{TZ},
provides the stability implication for nondegenerate positive
equilibria in full-rank three-species zero-one networks. That
criterion is prior work. Neither the structural observation nor
the bound in Theorem~\ref{thm:bound} is a new independent
certification of the reported bistability classification of all
373 remaining networks.

\section{One seven-reaction family}\label{sec:seven}
Add the reaction $X_3\xrightarrow{h}0$ to network 33449. With
the six original rates denoted $a,b,c,d,e,f$, the equations are
\begin{equation}\label{eq:seven}
\begin{aligned}
 \dot x&=-dx+by+az,\\
 \dot y&=-by-cyz+exz,\\
 \dot z&=-(a+h)z+by-cyz+fxy.
\end{aligned}
\end{equation}
All seven rates are strictly positive. This remains a full-rank,
three-species quadratic zero-one network.

\begin{proposition}\label{prop:seven}
For every $(a,b,c,d,e,f,h)\in\Rp^7$, system~\eqref{eq:seven}
has at most three distinct positive equilibria and at most one
exponentially stable positive equilibrium.
\end{proposition}
\begin{proof}
Put $t=z/y$, $L(t)=ea t+eb-cd$, and $D(t)=tL(t)$.
The first two equilibrium equations give the complete positive
branch
\begin{equation}\label{eq:branch7}
 t>\max\{0,(cd-eb)/(ea)\},\qquad
 x=\frac{b(b+at)}D,\quad y=\frac{bd}D,\quad z=\frac{bd}L.
\end{equation}
Indeed, the first equation gives $x=y(b+at)/d$, and substitution
in the second gives $y=bd/(tL)$. Positivity excludes $tL=0$
and is equivalent to the stated interval. On this branch the
third equation is $bd Q_h(t)/D(t)^2$, where
\begin{equation}\label{eq:q7}
\begin{split}
 Q_h(t)={}&-ea(a+h)t^3+[acd-h(eb-cd)]t^2\\
          &+b(eb+fa-2cd)t+b^2f.
\end{split}
\end{equation}
Its leading coefficient is strictly negative. Positive equilibria
are therefore in one-to-one correspondence with its roots in
the interval~\eqref{eq:branch7}, and there are at most three.

The coordinate change $z=ty$ has determinant $y$. After this
substitution, with $t$ fixed, the $(x,y)$ minor of the first two
equations is $-yD$ along their zero set. Differentiating
the third equation along the branch and evaluating at $Q_h=0$
gives
\begin{equation}\label{eq:det7}
 \det J=-yQ_h'(t).
\end{equation}
Exponential stability in three dimensions requires $\det J<0$,
so it requires $Q_h'>0$. The derivative is a downward-opening
quadratic; the set on which it is positive is empty or an interval, and
$Q_h$ is strictly increasing there. At most one root can lie in
that interval. A root with $Q_h'=0$ has singular Jacobian and
is not exponentially stable. This upper bound uses only a
necessary stability condition, not the sufficient direction of
Tang--Zhang's lemma.
\end{proof}

The independent derivation retains $z$ instead. With
$E(z)=db+(dc-be)z$, the original equations give
$ezF_1+dF_2=ae z^2-Ey$, so $E>0$ at every positive equilibrium.
They yield
\[
 x=\frac{az(b+cz)}E,\quad y=\frac{ae z^2}E,\quad
 P(z)=(a+h)E^2-aez(b-cz)E-a^2efz^2(b+cz).
\]
The remaining equation is $-zP/E^2$ and at an equilibrium
$\det J=-zP'/E$. Since $P(0)=(a+h)d^2b^2>0$ and $\deg P\leq3$,
there can be at most one positive root with $P'>0$: two upward
crossings would require a downward crossing before the first
and another between them, hence at least four distinct roots.
This gives the same stability bound through a different coordinate.

\subsection{An exact rational example}
Take
\begin{equation}\label{eq:rates}
 (a,b,c,d,e,f,h)=
 \left(\frac{109}{32768},\frac{235}{256},\frac12,\frac12,
       \frac12,\frac12,\frac1{65536}\right).
\end{equation}
Here the admissible interval is $t>0$, and
\[
 Q_h(t)=-\frac{23871}{4294967296}t^3
        +\frac{27797}{33554432}t^2
        -\frac{606065}{16777216}t+\frac{55225}{131072}.
\]
Exact Sturm sequences isolate one simple root in each interval
of Table~\ref{tab:roots}; the total number of admissible roots
is three. Rational interval evaluation verifies positive
concentrations and the displayed Hurwitz signs, both from
the symbolic branch and directly from the original Jacobian.
For $\det(\lambda I-J)=\lambda^3+b_1\lambda^2+b_2\lambda+b_3$,
positivity of $b_1,b_2,b_3,b_1b_2-b_3$ proves exponential stability.
If $b_3<0$, the characteristic polynomial has a positive real root.

\begin{table}[ht]
\centering\small
\begin{tabular}{@{}cllc@{}}
\toprule
Root & Rational lower endpoint & Rational upper endpoint &
$(b_1,b_2,b_3,b_1b_2-b_3)$\\
\midrule
1 & $1186617748/63717463$ & $1307453375/70205938$ & $(+,+,-,+)$\\
2 & $1998238606/38640527$ & $3149229703/60897580$ & $(+,+,+,+)$\\
3 & $6820268592/86644919$ & $6713775719/85292030$ & $(+,+,-,+)$\\
\bottomrule
\end{tabular}
\caption{Certified disjoint intervals in increasing $t$ order.
Exactly the middle equilibrium is exponentially stable.}
\label{tab:roots}
\end{table}

The independent $z$ intervals map to these roots in reverse order
under $t=db/(ae z)+(dc-be)/(ae)$. The exact transformed intervals
overlap in intervals containing one root each, and all four
Hurwitz signs agree. Thus the comparison checks corresponding
roots as well as their counts.

Persistence of nondegenerate equilibria under a sufficiently
small dependent reaction addition is consistent with the
inheritance results of Banaji--Pantea~\cite{BP}. The network has
full rank, so the added reaction vector is dependent on the
original stoichiometric vectors. The family-specific assertion
here is the upper bound for \emph{all} positive rates, together
with the exact input~\eqref{eq:rates}. We make no claim of a new
equivalence class or a classification of seven-reaction networks.

\section{Verification boundary and limitations}\label{sec:scope}
The independent implementation received the fixed upstream data,
definitions and statements to test. Its derivations, code and
results were fixed before comparison with the construction
records. The anticipated conclusions were disclosed, so this was
not a blind test of the hypotheses. The independent work used
different elimination and a separate implementation; it shared
the mathematical definitions, dataset and part of the SymPy
arithmetic infrastructure. Roles, dates, hashes and execution
records are in the companion \texttt{VERIFICATION-RECORD.md}.

\begin{table}[ht]
\centering\small
\begin{tabularx}{\linewidth}{@{}>{\raggedright\arraybackslash}p{30mm}YY@{}}
\toprule
Claim & Mathematical bridge & Finite verification\\
\midrule
375-network bound & Nonvanishing denominator, injective
projection, root bound & Cramer identities and signs;
independent resultants and Sylvester expansions for every input\\
Structural observation & Specified graph and finite domain &
All 375 graphs checked; precisely two selected\\
Seven-reaction bound & Complete positive branch and
Jacobian/derivative identity & Independent symbolic identity checks\\
Rational example & Sturm root counting and Hurwitz criterion &
Exact isolation, interval signs and root correspondence\\
\bottomrule
\end{tabularx}
\caption{What supports each assertion.}
\end{table}

The original enumeration yielding the list of 375 and the full
existence classification were not independently repeated. Our
bound applies exactly to that fixed list and does not by itself
prove the same assertion for every three-species, six-reaction
quadratic zero-one network. The 373-network bistability
classification is outside this independent certification.
The structural condition was selected after a pilot, and its
extension beyond the list is unresolved. Degenerate nonlinear
stability, boundary attractors, periodic orbits, dissipativity
and biological realization are not asserted. No human peer
review or Lean formalization has been completed.

\paragraph{Use of AI.}
Codex and Claude assisted derivations, implementation, verification,
literature comparison and drafting. Tatsuhiko Takimoto is responsible
for the manuscript. The assurances are the mathematical arguments
and exact checks specified above.

\begin{thebibliography}{9}\raggedright
\bibitem{TZ}
Xiaoxian Tang and Jiandong Zhang.
\emph{The Ubiquity of Three Steady States: Minimal Multistable
Zero-One Reaction Networks}. arXiv:2608.01116v1, 2026.
\url{https://arxiv.org/abs/2608.01116v1}.
\bibitem{TF}
Ang\'elica Torres and Elisenda Feliu.
\emph{Symbolic proof of bistability in reaction networks}.
arXiv:1909.13608v2, 2020.
\url{https://arxiv.org/abs/1909.13608v2}.
\bibitem{BF}
Murad Banaji and Elisenda Feliu.
\emph{Positive equilibria in mass action networks: geometry and bounds}.
Journal of Mathematical Biology 93, article 37, 2026.
Accepted manuscript: arXiv:2409.06877v5.
\url{https://arxiv.org/abs/2409.06877v5}.
\bibitem{BP}
Murad Banaji and Casian Pantea.
\emph{The inheritance of nondegenerate multistationarity in chemical
reaction networks}. arXiv:1608.08400v3, 2017.
\url{https://arxiv.org/abs/1608.08400v3}.
\bibitem{Adata}
Tatsuhiko Takimoto.
\emph{Stability certificates for two exceptional reaction networks}.
Prepared data-and-tools collection, version 1.0.0, 2026.
Repository name: \texttt{exception-networks-stability-certificates}.
Provisional reference: release and DOI pending
(\texttt{A-DOI-PENDING}; replace before publication).
\end{thebibliography}

\appendix
\section{Fixed input and certificate index}\label{app:input}
The upstream input is
\texttt{results/375\_multistable\_networks.json} in
\url{https://github.com/zjdong-sudo/paper-multistability-results},
at revision
\begin{center}\small
\path{323360df12c0787f9bbf7f223a6229b98233be39}.
\end{center}
Its SHA-256 is
\begin{center}\small
\path{517d1f16f890159185f5384fa2ab089ec4d270e03efe31161ed29489121e252c}.
\end{center}
The companion packet's \path{manifest.json} records hashes of its
files. The 375 candidate identifiers occur in both
\path{certificates/cramer.jsonl} and
\path{certificates/resultants.json}. Paths are relative to the
packet root. The latter file combines the independently obtained
factors and coefficient-sign records, including the three
equation-transform cases. The \path{certificates/} directory also
contains the structural table and rational-example certificates.
Run \path{tools/fetch_input.py} and then \path{verify.py}; the README
gives the environment and commands. The authors' raw data and code
are not bundled. The fetch tool obtains the identified input from
upstream and checks its SHA-256 before use.

This v1 is a local draft. Before public circulation, the certificate
packet and its inventory need a stable versioned public location;
no such archive or identifier is asserted here.
\end{document}
